\section{Proofs of lower and adaptation results}
\label{app:lower-proofs}

\begin{proof}[Proof of Theorem~\ref{thm:lower}]
Let \(\omega\in\{-1,1\}^d\) and set \(u_\omega=a\omega\), where
\[
a^2=c_0\left(\frac1n\wedge\frac{\rho^2}{d}\right)
\]
and \(c_0\) is sufficiently small. Then \(u_\omega\in B_d(0,\rho)\). If \(\omega\) and \(\omega^{(j)}\) differ only in coordinate \(j\), \cref{eq:kl-upper} gives
\[
\KL\{(Q_{u_\omega}^w)^{\otimes n}\|
        (Q_{u_{\omega^{(j)}}}^w)^{\otimes n}\}
\le 4C_Kna^2\le4C_Kc_0.
\]
Choose \(c_0\) so that Pinsker's inequality bounds the total variation between every neighboring product experiment by at most \(1/2\). Assouad's lemma then forces expected Hamming error at least \(d/4\) for any decoder of \(\omega\).

Given an arbitrary estimated conditional distribution \(\wt p\), let
\(\wh\omega\) be a Hellinger-nearest vertex, with measurable tie breaking.
Because \(p_{u_\omega}\) is one of the vertices, the triangle inequality gives
\[
h_w(p_{u_\omega},p_{u_{\wh\omega}})
\le 2h_w(p_{u_\omega},\wt p).
\]
Together with \cref{eq:hell-lower}, this implies
\[
h_w^2(p_{u_\omega},\wt p)
\ge c_h a^2d_H(\omega,\wh\omega),
\]
where \(d_H\) is Hamming distance. Therefore
\[
\sup_\omega\E_\omega h_w^2(p_{u_\omega},\wt p)
\ge c d a^2
=c\left(\rho^2\wedge\frac dn\right).
\]
Finally \(\Kw{q}{p}\ge h_w^2(q,p)\) under the convention
\(h^2=\int(\sqrt q-\sqrt p)^2\), which proves the claim. See~\citet{tsybakov2009} for the metric form of Assouad's lemma.
\end{proof}

\begin{proof}[Proof of Proposition~\ref{prop:normal-tube}]
To verify the asserted local existence of the normal field, choose a nonzero
vector \(v_0\) in the \(J_w\{\iota(0)\}\)-orthogonal complement of
\(D\iota(0)\). The matrices
\[
\Pi_u=I_D-D\iota(u)
\{D\iota(u)^\top J_w(\iota(u))D\iota(u)\}^{-1}
D\iota(u)^\top J_w(\iota(u))
\]
project onto the corresponding information-normal spaces. For the sufficiently
small \(\rho_0\) fixed in the proposition, \(\Pi_uv_0\) stays nonzero;
normalizing it in the \(J_w\)-metric
produces a \(C^2\) field satisfying \cref{eq:fisher-normal}. Compactness of
the smaller closed ball supplies the derivative bounds, and openness of
\(\mathcal U\) ensures \(\iota(u)+aN(u)\in\mathcal U\) for small \(a\).

The common-support and dominated-derivative assumptions give the uniform
two-parameter KL expansion
\begin{equation}
\Kw{p_\zeta}{p_{\zeta'}}
=\tfrac12(\zeta'-\zeta)^\top
J_w(\zeta_0)(\zeta'-\zeta)
+R(\zeta,\zeta';\zeta_0),
\qquad
|R|\le C(\norm{\zeta-\zeta_0}+\norm{\zeta'-\zeta_0})^3,
\label{eq:two-parameter-kl-expansion}
\end{equation}
whenever all three parameters lie in a sufficiently small common
neighborhood. This follows by differentiating KL in both arguments at
\((\zeta_0,\zeta_0)\): both first derivatives vanish and the second-order
quadratic form is the information applied to the parameter difference.

Take \(\zeta_0=\iota(u)\), \(\zeta=\zeta_{u,a}\), and
\(\zeta'=\iota(v)\). Taylor expansion of the embedding gives
\[
\zeta'-\zeta
=D\iota(u)(v-u)-aN(u)+O(\norm{v-u}^2).
\]
The cross term in the leading quadratic vanishes by
\cref{eq:fisher-normal}; information nondegeneracy and the singular-value
bounds for \(D\iota\) make the remaining quadratic uniformly equivalent to
\(a^2+\norm{u-v}^2\). The initial choice of \(\rho_0\), followed by shrinking
\(a_0\), absorbs the cubic
remainder and proves both sides of \cref{eq:tube-loss}.

The map \(u\mapsto\zeta_{u,a}\) has a uniformly bounded derivative. Applying
the upper side of \cref{eq:two-parameter-kl-expansion} to
\(p_{\zeta_{u,a}}\) and \(p_{\zeta_{v,a}}\) therefore gives
\[
\Kw{P_{u,a}}{P_{v,a}}\le C_T\norm{u-v}^2.
\]
The environment--prompt marginal is common, so this conditional KL equals
the one-observation joint KL; product additivity proves
\cref{eq:tube-kl}.

For the approximation statement, the candidate \(p_{\iota(u)}\) and the
upper side of \cref{eq:tube-loss} give
\(A_w^2(\lambda;P_{u,a})\le C_Ta^2\). The lower side gives
\(c_Ta^2\) for every candidate in the local chart. Shrink \(a_0\) further so
that \(c_Ta_0^2\le g_0\);
\cref{eq:off-chart-gap} excludes every remaining candidate and proves the
lower side of \cref{eq:tube-approx}.
\end{proof}

\begin{proof}[Proof of Theorem~\ref{thm:tube-lower}]
Put a hypercube \(u_\omega=\delta\omega\) in \(B_d(0,\rho)\), where
\[
\delta^2=c_0\left(\frac1n\wedge\frac{\rho^2}{d}\right).
\]
By \cref{eq:tube-kl}, neighboring product experiments have KL at most a
fixed constant when \(c_0\) is sufficiently small. Hence Assouad's lemma
forces Hamming error of order \(d\).

If \(\wt W\) lies in the local chart, write
\(\wt W=\iota(\wh v)\). If it lies outside, set \(\wh v=0\). In the first
case \cref{eq:tube-loss} gives
\[
\Kw{P_{u_\omega,a}}{p_{\wt W}}
\ge c_Ta^2+c_T\norm{u_\omega-\wh v}^2.
\]
In the second case, \cref{eq:off-chart-gap}, together with the choice of the
small local radii, bounds the same loss below by a sufficiently small constant
multiple of \(a^2+\norm{u_\omega-\wh v}^2\). Decode a nearest hypercube
vertex from \(\wh v\). The Euclidean version of the nearest-vertex argument
in the proof of Theorem~\ref{thm:lower} and Assouad's lemma give
\[
\sup_\omega\E_{u_\omega,a}
\Kw{P_{u_\omega,a}}{p_{\wt W}}
\ge ca^2+cd\delta^2.
\]
Since the hypercube is a subset of the ball, this is
\cref{eq:tube-lower-result}.
\end{proof}

\begin{proof}[Proof of Proposition~\ref{cor:scale}]
For a fixed integer \(D\), the truncation estimator
\(\wh\theta_j=Z_j\mathbf 1\{j\le D\}\) has risk
\[
\frac Dn+\sum_{j>D}\theta_j^2
\le \frac Dn+L^2D^{-2s}.
\]
Taking the best \(D\) proves the upper bound.

For the lower bound, first suppose the variance--bias crossing occurs at an
integer \(D\ge1\), up to a constant factor depending only on \(s\):
\[
D/n\asymp L^2D^{-2s}.
\]
Place a hypercube on the shell \(j=D+1,\ldots,2D\), with each active
coordinate equal to \(\pm\delta\), where
\[
\delta^2=c_s\min\left\{\frac1n,
L^2D^{-(2s+1)}\right\}.
\]
For sufficiently small \(c_s\), every vertex belongs to
\(\mathcal H_s(L)\): if \(m\le D\), its tail energy is at most
\(D\delta^2\le L^2D^{-2s}\le L^2m^{-2s}\); if
\(D<m<2D\), use \(m^{-2s}\ge(2D)^{-2s}\) and decrease \(c_s\); for
\(m\ge2D\) the tail is zero. Neighboring Gaussian experiments have KL
\(2n\delta^2\), which is bounded by a sufficiently small constant.
Assouad's lemma therefore yields risk at least
\[
cD\delta^2
\asymp c_s\left\{D/n+L^2D^{-2s}\right\},
\]
where the final equivalence uses the crossing relation. Monotonicity of the
two terms shows that this is within a constant factor of their infimum.

If the continuous crossing lies below one, then \(L^2<1/n\) and the displayed
infimum is comparable to \(1/n\). A two-point experiment in the unconstrained
first coordinate, at separation \(c/\sqrt n\), gives this lower bound.
Rounding the crossing to an integer changes constants only. The bounded-ratio
grid statement follows by rounding \(D\) upward to the next grid point.
\end{proof}

\begin{proof}[Proof of Theorem~\ref{thm:lan}]
Differentiability in quadratic mean and nonsingular \(I_\pi\) imply convergence of the local experiments to the Gaussian shift
\[
Z=h+I_\pi^{-1/2}\xi,\qquad \xi\sim N(0,I_d).
\]
For every \(M<\infty\), apply the local asymptotic minimax theorem to the
bounded bowl-shaped loss
\[
\ell_M(a,h)=\min\{M,\tfrac12(a-h)^\top G_w(a-h)\}.
\]
The assumed uniform truncated-loss convergence and coercivity justify the
transfer from the original estimator range to the local Gaussian experiment
\citep{vandervaart1998}. Letting the local ball expand and then
\(M\to\infty\), monotone convergence gives the minimax risk of the untruncated
Gaussian shift. The equivariant estimator \(\wh h=Z\) has risk
\[
\tfrac12\E\xi^\top I_\pi^{-1/2}G_wI_\pi^{-1/2}\xi
=\tfrac12\tr(G_wI_\pi^{-1}),
\]
and is minimax in the expanding-ball limit. An asymptotically linear efficient
estimator, uniformly on compact local alternatives, attains the same constant
when its rescaled losses are uniformly integrable. At correct conditional
specification, the target Hessian is \(I_w\); setting \(w=\pi\) gives
\(\tr(I_d)/2=d/2\).
\end{proof}

\begin{proof}[Proof of Proposition~\ref{prop:misspecified-constant}]
The asserted asymptotic linearity and central limit theorem give covariance
\[
\Sigma_\pi=H_\pi^{-1}V_\pi H_\pi^{-1}
\]
for \(\sqrt n(\wh\theta-\theta^*)\). Because \(\theta^*\) is stationary, a
second-order expansion of the target risk gives
\[
R_w(\wh\theta)-R_w(\theta^*)
=\tfrac12(\wh\theta-\theta^*)^\top
G_w(\wh\theta-\theta^*)+o_p(n^{-1}).
\]
Uniform integrability permits expectations to pass through this expansion.
The Gaussian quadratic-form identity and cyclicity of trace then yield
\[
\lim_{n\to\infty}n\E\{R_w(\wh\theta)-R_w(\theta^*)\}
=\tfrac12\tr(G_w\Sigma_\pi)
=\tfrac12\tr\{G_wH_\pi^{-1}V_\pi H_\pi^{-1}\}.
\]
\end{proof}

\begin{proof}[Proof of Theorem~\ref{thm:rank-impossibility}]
Any rank selector induces a test between \(W_0\) and \(W_t\): accept the higher-rank alternative when it returns \(\bm r\). The larger of its two rank errors is at least half the sum of the two testing errors. The optimal sum is
\(1-\TV\{(Q_{W_0}^\pi)^{\otimes n},(Q_{W_t}^\pi)^{\otimes n}\}\). By product additivity of KL and Pinsker's inequality,
\[
\TV\{(Q_{W_0}^\pi)^{\otimes n},(Q_{W_t}^\pi)^{\otimes n}\}
\le\sqrt{\frac n2\KL(Q_{W_0}^\pi\|Q_{W_t}^\pi)}
\le\sqrt{Cn\delta_t^2/2}.
\]
Substitution proves the result.
\end{proof}

\begin{proof}[Proof of Proposition~\ref{prop:rank-sufficiency}]
The criterion of the true rank is at most
\(P_n\ell_W+\operatorname{pen}_{\bm r}/n\). On the asserted event, every
underfit criterion exceeds this upper bound by at least
\[
g_n/2+(\operatorname{pen}_{\bm s}-
\operatorname{pen}_{\bm r})/n>0,
\]
and every overfit criterion exceeds it by at least
\[
(\operatorname{pen}_{\bm s}-\operatorname{pen}_{\bm r}
-b_{\bm s,n})/n>0.
\]
Thus \(\wh{\bm r}=\bm r\) on events whose probabilities tend to one.
\end{proof}
